\documentclass[11pt]{article}

\usepackage[margin=1in]{geometry}
\usepackage{booktabs}
\usepackage{amsmath}
\usepackage{array}
\usepackage{hyperref}

\title{CS234 Results Draft}
\date{2026-03-15}

\begin{document}
\maketitle

\section{Results}

We ran the full-scale pipeline on the complete multiple-choice quiz bowl dataset
produced by our preprocessing pipeline. After answer-profile construction and
MC filtering, the final dataset contained \textbf{14,961} tossups, split into
\textbf{10,453 train / 2,226 validation / 2,282 test}. For the belief-feature
pipeline, all reported baseline and PPO experiments used \textbf{TF-IDF}
likelihoods; for the end-to-end text policy pipeline, we trained a
\textbf{T5-base} model on Apple Silicon MPS.

\subsection{Baseline Performance}

Our strongest classical baseline was the \textbf{SequentialBayesBuzzer}, which
achieved the best overall quiz-bowl utility score of
\textbf{$S_q = 0.7413$} at threshold 0.5, with \textbf{96.97\%} buzz accuracy,
\textbf{ECE $= 0.3470$}, and \textbf{Brier $= 0.1535$}. By comparison, both the
\textbf{ThresholdBuzzer} and \textbf{SoftmaxProfileBuzzer} peaked at
\textbf{$S_q = 0.6858$} with \textbf{97.48\%} buzz accuracy. This gap suggests
that belief tracking matters more than simply thresholding top-answer
confidence. In particular, Sequential Bayes trades a small drop in raw accuracy
for substantially better timing behavior, which is exactly what $S_q$ is
intended to reward.

\begin{table}[h]
\centering
\begin{tabular}{lrrrrr}
\toprule
Agent & Threshold & Accuracy & $S_q$ & ECE & Brier \\
\midrule
ThresholdBuzzer & 0.5 & 0.9748 & 0.6858 & 0.4555 & 0.2326 \\
SoftmaxProfileBuzzer & 0.5 & 0.9748 & 0.6858 & 0.4555 & 0.2326 \\
SequentialBayesBuzzer & 0.5 & 0.9697 & 0.7413 & 0.3470 & 0.1535 \\
\bottomrule
\end{tabular}
\caption{Best baseline results on the TF-IDF belief-feature pipeline.}
\end{table}

\noindent\textbf{Takeaway.} Among belief-feature methods, Sequential Bayes is
the strongest overall policy under our utility metric.

\subsection{PPO Performance}

In contrast, the default PPO agent performed poorly. Its final test performance
was only \textbf{25.25\%} accuracy with \textbf{$S_q = 0.2525$} and negative
reward-like return (\textbf{$-0.2777$}). This is far below all three baseline
families and indicates that, under our current state representation and
training budget, PPO is not learning a useful buzzing policy. The low ECE
(\textbf{0.0025}) is misleading here: the model is not well calibrated in a
meaningful decision-theoretic sense, but rather collapsed into a low-confidence
regime that produces a numerically small calibration gap.

\begin{table}[h]
\centering
\begin{tabular}{lrrrr}
\toprule
Policy & Accuracy & $S_q$ & Reward-like & ECE \\
\midrule
PPO default & 0.2525 & 0.2525 & -0.2777 & 0.0025 \\
Best baseline (SequentialBayes) & 0.9697 & 0.7413 & --- & 0.3470 \\
\bottomrule
\end{tabular}
\caption{Default PPO versus the best TF-IDF baseline.}
\end{table}

\noindent\textbf{Takeaway.} PPO, as implemented here, is not competitive with
the engineered baseline pipeline.

\subsection{End-to-End T5 Policy}

The end-to-end T5 policy tells a more nuanced story. After supervised warm-start
and PPO fine-tuning, the \textbf{T5 policy} reached \textbf{93.34\%} answer
accuracy, far above the MLP policy used in the belief-feature RL pipeline
(\textbf{25.46\%}). However, its quiz-bowl utility remained low, with
\textbf{$S_q = 0.2394$}, \textbf{ECE $= 0.6816$}, and
\textbf{Brier $= 0.5268$}. The model appears to be much better at eventually
producing the correct answer, but much worse at calibrated buzzing and timing.
By comparison, the MLP policy had lower accuracy but slightly higher $S_q$
(\textbf{0.2546}) simply because both policies are weak under the
timing-sensitive objective.

\begin{table}[h]
\centering
\begin{tabular}{lrrrrr}
\toprule
Model & Accuracy & $S_q$ & ECE & Brier & Avg buzz pos \\
\midrule
MLP policy (T5-as-likelihood) & 0.2546 & 0.2546 & 0.0046 & 0.1898 & 0.0000 \\
T5 policy (end-to-end) & 0.9334 & 0.2394 & 0.6816 & 0.5268 & 3.8966 \\
\bottomrule
\end{tabular}
\caption{Phase 6 comparison between the belief-feature MLP policy and the
end-to-end T5 policy.}
\end{table}

\noindent\textbf{Takeaway.} The end-to-end text policy is strong at final
correctness but not yet strong at calibrated early buzzing.

\subsection{Expected Wins}

Against a logistic opponent model, our Expected Wins evaluation gave a mean
score of \textbf{4.4929}. This experiment is best interpreted as an alternate
downstream metric rather than a new training success: it evaluates the same
baseline family under a different reward model. We did not complete the
separate EW-trained PPO phase in this run, so these results should be read as
evaluation-only, not as evidence that PPO improved under Expected Wins
training.

\subsection{Ablation Studies}

Our ablations reinforce the central result that baseline design matters more
than PPO configuration.

\paragraph{Reward-mode ablations.}
Switching PPO from the default time-penalty reward to either \texttt{simple} or
\texttt{human\_grounded} produced \textbf{no measurable improvement}: all three
variants stayed at roughly \textbf{25.25\%} accuracy and
\textbf{$S_q = 0.2525$}. Under the current setup, reward reshaping alone is not
enough to recover a strong learned policy.

\paragraph{Belief-mode ablation.}
Replacing the default from-scratch confidence with \textbf{sequential Bayes
belief updates} increased $S_q$ from \textbf{0.6858} to \textbf{0.7413},
making this the most effective single ablation in the entire run.

\paragraph{Stop-only PPO.}
The stop-only policy produced a striking failure mode: \textbf{97.55\%}
accuracy but only \textbf{$S_q \approx 2.9 \times 10^{-5}$}. In other words,
the agent learned behavior that looks good under eventual correctness but is
essentially useless under the quiz-bowl timing metric. This is a clear example
of reward/objective mismatch.

\paragraph{No-buzz horizon.}
Allowing a no-buzz end state did not help. The no-buzz variant remained
effectively identical to the default PPO model, with \textbf{25.15\%} accuracy
and \textbf{$S_q = 0.2515$}.

\subsection{Variance Across Seeds}

PPO training was highly unstable across seeds. Two seeds converged to the
low-performing default regime (\textbf{24.45\%} accuracy,
\textbf{$S_q \approx 0.2445$}), while one seed converged to the same
high-accuracy / near-zero-$S_q$ failure mode seen in the stop-only variant
(\textbf{97.55\%} accuracy, \textbf{$S_q = 0.000069$}). Across seeds, the mean
$S_q$ was \textbf{0.1630} with population standard deviation \textbf{0.1152} ($n=3$), confirming
that PPO behavior is not only weak on average but also qualitatively
inconsistent.

\noindent\textbf{Takeaway.} The learned RL policy is unstable and often
converges to degenerate local optima.

\subsection{Distractor and K-Sensitivity Analyses}

The distractor study produced a somewhat surprising result. With the same
Sequential Bayes baseline, \textbf{category-random distractors} gave the best
score (\textbf{$S_q = 0.7962$}), \textbf{SBERT-ranked distractors} matched the
default setup (\textbf{0.7413}), and \textbf{TF-IDF-profile distractors} were
the hardest (\textbf{0.6651}). This suggests that the exact negative-answer
construction strategy materially changes task difficulty.

We also tested fixed and mixed values of $K$, the number of answer choices. On
the mixed-$K$ dataset, Sequential Bayes reached \textbf{$S_q = 0.7804$}.
Across fixed $K$, performance improved over the default $K=4$ setting and
peaked around \textbf{$K=5$} with \textbf{$S_q = 0.8000$}, then remained
similarly high at \textbf{$K=6$} (\textbf{0.7983}). This suggests that the
default $K=4$ setting is not necessarily the easiest or best-conditioned
version of the task.

\subsection{Overall Interpretation}

The overall picture is clear. The \textbf{belief-feature baseline pipeline},
especially \textbf{Sequential Bayes}, is currently the strongest approach in
this project. PPO-based learned policies are either weak, unstable, or
vulnerable to objective mismatch. The end-to-end T5 policy demonstrates that
high answer accuracy is achievable, but it does not yet translate into good
buzzing utility. For this project, the most defensible conclusion is that
\textbf{carefully engineered uncertainty and belief-tracking methods outperform
our current RL formulations on the quiz-bowl buzzing objective}.

\section{Discussion and Limitations}

Two implementation issues surfaced during the live run and were fixed: an Apple
Silicon MPS device-placement bug in Phase 6 and a \texttt{StopOnlyEnv}
evaluation bug in Phase 16. These did not change the final substantive
conclusions, but they do highlight that full-scale execution exposed issues
that were not caught by smoke tests alone.

The main empirical limitation is that our PPO results are based on a relatively
modest training budget and a state/action design that may be poorly matched to
the true structure of the task. In addition, some optional extensions,
including EW-trained PPO and DSPy/OpenAI-based paths, were not run in this
full-scale pass. Future work should focus on stabilizing policy learning,
aligning the reward more tightly with $S_q$, and improving calibration for the
end-to-end T5 model.

\end{document}
